\chapter{法术、精灵与职业技能命中机制}

物理系统和法术系统在当前源码里共用“命中/未命中”这一结果，但它们进入结果的变量层次并不相同。物理判定主要围绕固定敏和固定幸运展开，而法术系统则大量改用原始幸运、属性抗性、抗魔装备和技能熟练度。本章的目标正是把这两套命中逻辑拆开说明，并进一步解释“法术命中”和“法术伤害”为什么必须分层处理。

\section{普通属性法术命中}

普通属性法术与精灵伤害的命中入口可统一理解为 \code{BATTLE_MagicDodge()}。若目标是玩家，则未命中基数写成
\begin{equation}
    m_{\mathrm{normal}}
    = \left\lfloor 3L + 0.15R + 0.9E \right\rfloor,
    \label{eq:magic-miss-player}
\end{equation}
其中 $L$ 是目标原始幸运 \code{CHAR_LUCK}，$R$ 是目标对应属性抗性，$E$ 是抗魔装备值。法术命中率于是为
\begin{equation}
    p_{\mathrm{magic}}^{(\mathrm{player})}=100-m_{\mathrm{normal}}.
    \label{eq:magic-hit-player}
\end{equation}

若目标不是玩家，则系统不再读取幸运、抗性和抗魔装备，而改用等级近似：
\begin{equation}
    m_{\mathrm{normal}}^{(\mathrm{npc})} = \min\left\{\left\lfloor 0.2L_v \right\rfloor,\,30\right\},
\end{equation}
从而
\begin{equation}
    p_{\mathrm{magic}}^{(\mathrm{npc})}=100-m_{\mathrm{normal}}^{(\mathrm{npc})}.
\end{equation}
这两式说明了一个经常被玩家忽略的事实：普通法术打人和打宠、打怪，并不是“同一公式只换目标名字”，而是直接换了输入变量集。对玩家来说，幸运、抗性、抗魔装都重要；对非玩家单位来说，等级却成了主要近似项。还要特别指出的是，普通法术命中这一步完全不读取生命值：无论是当前生命 \code{CHAR_HP}、最大生命 \code{CHAR_WORKMAXHP}、生命成长档，还是“去掉装备修正后的生命口径”，都不进入 \code{BATTLE_MagicDodge()}。因此，生命高低本身不会改变普通法术是否命中；它只会在后续的伤害、治疗或少数以生命为基准的特殊技能里发挥作用。

\section{精灵伤害命中}

在当前实现中，精灵伤害和普通属性法术共用同一类命中入口，因此可以与式~\eqref{eq:magic-miss-player}--式~\eqref{eq:magic-hit-player} 同口径理解。也就是说，若目标是玩家，精灵是否命中主要仍由原始幸运、对应抗性和抗魔装备决定；若目标不是玩家，则主要退化为等级近似。

这意味着精灵伤害的“命中优化”空间其实很有限。进攻方通常很难像物理系那样主动堆一个明确的命中主属性，而更常见的是从守方的抗性、抗魔装和幸运入手分析减伤与抗性配置。因此，精灵命中与其说是进攻方的构筑问题，不如说更接近防守方的属性抗性问题。

\section{职业技能命中与熟练度}

职业法术命中改由职业法术闪避判定流程负责。与普通法术相比，它最重要的新增变量是对应属性熟练度。若目标是玩家，则第一段未命中基数可整理为
\begin{equation}
    m_{\mathrm{prof}}
    = \left\lfloor 3L + 0.5R + 0.4E - 0.2P \right\rfloor,
    \label{eq:prof-miss-player}
\end{equation}
其中 $L$ 为目标原始幸运，$R$ 为目标对应职业属性抗性，$E$ 为抗魔装备，$P$ 为攻方对应属性熟练度。于是第一段命中率为
\begin{equation}
    p_{\mathrm{prof},1}=100-m_{\mathrm{prof}}.
    \label{eq:prof-hit-stage1}
\end{equation}

若目标不是玩家，则幸运、装备项退化为等级近似，得到
\begin{equation}
    m_{\mathrm{prof}}^{(\mathrm{npc})}
    = \left\lfloor \min\{0.15L_v,20\} - 0.2P \right\rfloor.
\end{equation}
因此，职业法术相比普通法术多了一个清晰可控的进攻端变量 $P$。从边际结构看，玩家目标每增加 $1$ 点原始幸运会让第一段命中率约下降 $3\%$，而攻方每增加 $5$ 点熟练度会让第一段命中率约提高 $1\%$。这也是职业法术研究中“熟练度比幸运更像可优化变量”的根本原因。与普通法术一样，职业法术命中这一步同样完全不读取生命值：当前生命 \code{CHAR_HP}、最大生命 \code{CHAR_WORKMAXHP}、生命成长档，以及去掉装备修正后的生命口径，都不进入 \code{PROFESSION_MAGIC_DODGE()}；它们只会在后续伤害或特殊技能效果里发挥作用。

还必须注意一个绝对前置条件：若目标当前处于 地球一周，则职业法术会直接 Miss，不再进入式~\eqref{eq:prof-miss-player} 的概率计算。也就是说，某些防御动作并不是降低命中率，而是直接切断命中流程。

把两类法术命中并列起来，可以得到表~\ref{tab:magic-hit-compare}。

\begin{table}[H]
    \centering
    \caption{普通法术与职业法术命中的输入变量对照}
    \label{tab:magic-hit-compare}
    \begin{tabularx}{\textwidth}{>{\raggedright\arraybackslash}p{3.2cm} X X}
        \toprule
        项目 & 普通属性法术 / 精灵 & 职业法术 \\
        \midrule
        命中入口 & \code{BATTLE\_MagicDodge()} & \code{PROFESSION\_MAGIC\_DODGE()} \\
        玩家目标主要输入 & 原始幸运、对应属性抗性、抗魔装备 & 原始幸运、对应职业属性抗性、抗魔装备、攻方熟练度 \\
        非玩家目标主要输入 & 目标等级近似，最高压到 $30$ & 目标等级近似，最高压到 $20$，再减熟练度项 \\
        是否读取生命值 & 不读取 & 不读取 \\
        额外前置绝对 Miss & 无这一步 & 若目标当前动作是 \code{BATTLE\_COM\_S\_EARTHROUND0}，直接 Miss \\
        二段命中 & 无 & 部分技能存在固定二段命中 \\
        \bottomrule
    \end{tabularx}
\end{table}

\section{状态技能与二段判定}

职业技能里并不是所有法术在通过第一段命中后就直接生效。若把式~\eqref{eq:prof-hit-stage1} 的第一段命中率记为 $p_{\mathrm{prof},1}$，则部分技能还会叠加一个固定第二段命中率 $p_2$，于是最终命中率写成
\begin{equation}
    p_{\mathrm{final}} = p_{\mathrm{prof},1}\, p_2.
    \label{eq:prof-hit-stage2}
\end{equation}
当前已核对的典型二段判定包括：
\begin{equation}
    p_2=
    \begin{cases}
        0.75, & \text{电流术、暴风雨},\\
        0.90, & \text{火龙枪、世界末日},\\
        0.50, & \text{一针见血}.
    \end{cases}
\end{equation}

冰爆术则是一个非常特别的例子。源码在原本的命中分支之后又直接把 \code{hit} 重写成 $100$，因此就当前代码口径而言，这一段额外命中视为必定通过。这个例子说明，研究法术命中不能只靠技能说明文本，而必须回到实际实现顺序，因为后续覆写足以改变整段逻辑。

对状态技能而言，还要再区分“命中成功”和“状态参数如何写入”。例如四属性结界会在命中后写入结界槽位；挑拨、树根缠绕等则会继续改写攻防敏工作值。换言之，命中本身只是状态技能的入口，不是它的最终效果。

\section{命中与伤害的分层关系}

法术是否命中，与命中后造成多少伤害，是两层彼此独立的计算。若把法术基础威力记为 $M$，攻方对应熟练度记为 $P$，守方对应抗性、抗装、精灵抗性分别记为 $R,S,G$，则项目当前整理出的职业法术伤害层可写为
\begin{equation}
    damage = M\left(1+\frac{P}{100}\right)\left(1-\frac{R}{100}\right)\left(1-\frac{S}{100}\right)\left(1-\frac{G}{100}\right).
    \label{eq:magic-damage-layer}
\end{equation}
式~\eqref{eq:magic-damage-layer} 与前面的命中公式共同说明：提高法术表现有两条完全不同的路径。第一条是让 $p_{\mathrm{final}}$ 上升，也就是更容易打中；第二条是让式~\eqref{eq:magic-damage-layer} 上升，也就是打中了更疼。熟练度同时参与这两层，因此往往是职业法术里最具复合收益的进攻变量；而抗性则既能降低第一段命中，又能在命中后继续减伤，因此同样具有双重防守价值。

从研究方法看，把命中层和伤害层分开还有一个更深的意义：它让“法术最优化”不再只是一个单变量问题。若目标是稳定控场，那么提高第一段命中比堆最终伤害更重要；若目标是单次爆发，则可能宁可接受较低命中率，也要换取更高的式~\eqref{eq:magic-damage-layer}。这种目标差异在后文的优化章节中会再次出现。

